\newpage\SourceLeaf{42}\ArithmeticHead{35}{R}{KREISTEILUNG.}\HistoricalFoot{}{5*}
\begin{center}\bfseries 14.\end{center}
\noindent\hspace*{\parindent}\textsc{Gauss}’ Methode beruht nun einerseits auf der Irreduzibilität der Gleichung \(X=0\), für welche er in art. 341 der D. A. den ersten der vielen Beweise gab, die man jetzt dafür kennt; nach Nr. 40 des \textit{Tagebuchs} war er aber erst am 9. Oktober 1796 im Besitze dieses Beweises und hat dann (Nr. 136 T.) am 12. Juni 1808 auch für die Irreduzibilität der Gleichung, welche die primitiven \(p^{\mathrm{ten}}\) Wurzeln der Einheit bestimmt, falls \(p\) eine zusammengesetzte Zahl ist, einen Beweis gehabt. Diese Gleichung selbst sowie einige Grundgedanken des Beweises finden sich in dem im Oktober 1805 begonnenen und über 1808 hinaus fortgeführten Handbuche (18, Bd) des \textsc{Gauss}schen Nachlasses (siehe W. X 1, S. 116); man darf mithin wohl der betreffenden Notiz auch jenes Datum des Tagebuchs beilegen.

Der andere Grundpfeiler seiner Theorie ist die Einteilung der Wurzeln in sogenannte Perioden. Ist \(p-1=e\mathbin{.}f\), so erhält man \(e\) Perioden von \(f\) Gliedern:
\[
\begin{gathered}
\eta_i=r^{g^i}+r^{g^{i+e}}+r^{g^{i+2e}}+\ldots+r^{g^{i+(f-1)e}}\\
(i=0,1,2,\ldots,e-1).
\end{gathered}
\]
Sie sind die Wurzeln einer Gleichung \(X_e=0\) vom \(e^{\mathrm{ten}}\) Grade, deren Koeffizienten rationale ganze Zahlen sind, und welcher dieselbe charakteristische Grundeigenschaft zukommt, wie der Gleichung \(X=0\), nämlich dass in der zyklischen Anordnung \(\eta_0,\eta_1,\eta_2,\ldots,\eta_{e-1}\) ihrer Wurzeln jede die gleiche rationale Funktion von der vorhergehenden ist. Diese Gleichung gestattet somit eine entsprechende Behandlung. Zerlegt man die \(f\)-gliedrigen Perioden, indem man \(e=e'\mathbin{.}f'\) setzt, jede in \(e'\) Perioden von \(f'\) Gliedern, so sind die \(e'\) Perioden
\[
\begin{gathered}
\eta'_i=r^{g^i}+r^{g^{i+ee'}}+\ldots+r^{g^{i+(f'-1)ee'}},\\
(i=0,1,2,\ldots,e'-1)
\end{gathered}
\]
welche die Periode \(\eta_0\) zusammensetzen, die Wurzeln einer Gleichung \(X'_{e'}=0\) vom Grade \(e'\), deren Koeffizienten rational (genauer: linear) durch die \(f\)-gliedrigen Perioden ausdrückbar und somit nach Auflösung der Gleichung \(X_e=0\) bekannt sind. Auch diese Gleichung ist von demselben Charakter, wie die
\newpage\SourceLeaf{43}\ArithmeticHead{36}{L}{BACHMANN, ÜBER GAUSS’ ZAHLENTHEORETISCHE ARBEITEN.}
\noindent vorigen, und kann daher ebenso behandelt werden, u. s. w. Hiernach kommt, wenn irgendwie
\[
p-1=ee'e''\ldots e^{(\nu-1)}
\]
gesetzt wird, die Auflösung der Kreisteilungsgleichung darauf zurück, der Reihe nach die Hilfsgleichungen
\[
X_e=0,\ X'_{e'}=0,\ldots,X^{(\nu-1)}_{e^{(\nu-1)}}=0
\]
von den Graden \(e,e',\ldots,e^{(\nu-1)}\) aufzulösen. Man erhält die Hilfsgleichungen kleinsten Grades, wenn man \(p-1\) als Produkt von Primzahlpotenzen:
\[
p-1=a^\alpha b^\beta\ldots c^\gamma
\]
darstellt, und hat dann \(\alpha\) Gleichungen vom Grade \(a\), \(\beta\) Gleichungen vom Grade \(b,\ldots,\gamma\) Gleichungen vom Grade \(c\) zu lösen. Nur in dem Falle \(p=2^{2^k}+1\) sind alle Hilfsgleichungen quadratisch, die Kreisperipherie also mittels Zirkel und Lineal in \(p\) gleiche Teile zerlegbar; in jedem andern Falle tritt mindestens eine Hilfsgleichung höheren Grades auf, und mit gesperrtem Druck ist in art. 365 von \textsc{Gauss} bemerkt worden, dass sie dann auf keine Weise zu vermeiden oder zu erniedrigen ist (omni rigore demonstrare possumus, has aequationes elevatas nullo modo nec evitari nec ad inferiores reduci posse). Dieser Passus ist offenbar dem Werke noch in letzter Stunde während des Drucks eingefügt worden, da (nach Nr. 116 T.) \textsc{Gauss} erst am 6. April 1801 zu solcher Gewissheit gelangte, die D. A. aber im Sommer dieses Jahres erschienen, und er zeigt, wie tiefe Blicke \textsc{Gauss} schon in die algebraische Auflösung der Gleichungen getan haben muss. Der Beweis selbst ist aber weder in den Schriften noch im Nachlasse von \textsc{Gauss} vorhanden. Das Siebzehneck gehört zu den erstbezeichneten Fällen, ist also mit Zirkel und Lineal konstruierbar; es ist in Übereinstimmung mit der \textsc{Gauss}schen Theorie zuerst von \textsc{Paucker}, dann von \textsc{Erchinger} rein geometrisch konstruiert worden\footnote{\textsc{Magnus Georg v. Paucker}, \textit{Geometrische Verzeichnung des regelmässigen 17-Ecks und 257-Ecks in den Kreis}, Jahresverhandlungen der Kurländischen Gesellschaft für Litteratur und Kunst \textbf{2}, 1822, S. 160—219; \textit{Die ebene Geometrie der geraden Linie und des Kreises} I, Königsberg 1823, S. 187. Die Abhandlung von \textsc{Johannes Erchinger}, über die \textsc{Gauss} 1825 in den Göttingischen gelehrten Anzeigen berichtet hat (siehe W. II, S. 186), ist wie es scheint überhaupt nicht gedruckt worden. Diese Vermutung stützt sich auf eine Bemerkung, die der Tübinger Professor der Rechte \textsc{Heinrich Schrader} in einem an \textsc{Gauss} gerichteten Briefe vom 1. September 1825 macht. Der Brief handelt eingehend von \textsc{Erchingers} Lebensschicksalen und von seinem in Rede stehenden Aufsatz, von dem gesagt wird: »dass er ohne neue Ausarbeitung der äussern Form nach kaum druckreif sein mögte«. Als älteste Konstruktion des regelmässigen Siebzehnecks dürfte die aus dem Jahre 1802 stammende und W. X 1, S. 120 zum ersten Male veröffentlichte des Tübinger Professors der Mathematik \textsc{Christoph Friedrich v. Pfleiderer} anzusehen sein.}.
